\section{Introduction}
\label{sec:intro}

% target: ~1 page

% paragraph 1 — the three-property gap
CATE estimation on real outcome data demands three properties that
few tools jointly provide: heterogeneous $\tau(x)$, calibrated
uncertainty on $\tau(x)$, and robustness to heavy tails. Meta-learners
\citep{kunzel2019metalearners} target the first; bootstrap CIs from
causal forests \citep{wager2018estimation} target the second
approximately; classical Bayesian causal models
\citep{hill2011bayesian, hahn2020bcf} target the second exactly but
under Gaussian-tail assumptions that fail on medical outcome
distributions with rare severe events, revenue distributions with
long-tailed customer value, and response surfaces with structural
heterogeneity. Real outcome data almost always has real tails, and
Gaussian-likelihood Bayesian models systematically misfit them.

% paragraph 2 — our position
This paper introduces the \emph{Bayesian X-Learner}: an X-Learner
whose Phase-3 inference layer is a full MCMC posterior over
$\tau(x)$ with a Welsch-form redescending pseudo-likelihood
\citep{dennis1978techniques}. Compared to Gaussian Bayesian baselines, the
redescender trades a known efficiency cost on truly clean data
\citep{huber1964robust, huber2009robust} --- the asymptotic relative
efficiency for a univariate location estimand transfers, with
finite-sample amplification we document empirically, to the CATE
setting --- for non-trivial robustness to heavy-tailed
pseudo-outcomes. The resulting posterior
is interpretable as a generalised Bayesian (Gibbs) posterior and supports the
downstream tasks --- subgroup comparisons, integrated decisions,
propagated uncertainty --- that point estimates with bootstrap CIs
do not.

% paragraph 3 — two complementary heavy-tail regimes
We distinguish two heavy-tail regimes the practitioner may face.
\emph{Tails-as-contamination}: outcomes contain a small fraction of
units whose $Y$ is distorted by nuisance mechanisms
(measurement error, reporting artefacts, revenue whales in marketing)
and whose true $\tau$ is nominal. \emph{Tails-as-signal}: a minority
subgroup has a genuinely different $\tau$, and the practitioner cares
about recovering that subgroup's effect. These regimes look
superficially identical but call for opposite tools: the first wants
downweighting (Welsch + Huber nuisance); the second wants
representation (a tail-aware CATE basis). A controlled DGP probe
(Section~\ref{sec:experiments}) confirms both pathways work as
hypothesised, and that the library's data-layer
\texttt{normalize\_extremes} path is a contamination-reduction op that
should not be activated for heterogeneous $\tau$.

% paragraph 4 — contributions
Our contributions are:
\begin{enumerate}
  \item A Bayesian X-Learner that pairs doubly robust cross-fitted
    pseudo-outcomes \citep{kennedy2020optimal} with a Welsch MCMC
    posterior over $\tau(x)$. On Hill's IHDP benchmark, the default
    configuration leads on mean $\PEHE = 0.56$ against S-/T-/X-
    Learner, full-config Causal BART \citep{hill2011bayesian} at
    0.60, and EconML's causal forest at 1.06, at the tightest
    dispersion among competitive entries.
  \item An explicit separation of \emph{tails-as-contamination} and
    \emph{tails-as-signal} regimes, with a controlled probe
    (Section~\ref{sec:experiments:tail_signal}) and a graded
    basis-sensitivity ablation
    (Section~\ref{sec:experiments:basis_ablation}) showing which
    library mechanism addresses which and why a data-layer
    Hill-estimator rescaling is actively harmful for the signal
    regime.
  \item Empirical validation of Huber's 1964 minimax-$\delta$
    prescription as an operational configuration: recovery at up to
    $\sim$20--25\% whale density via a nuisance-loss selection
    exposed through a \texttt{contamination\_severity} enum
    (Section~\ref{sec:method:severity}), with documented breakdown
    beyond 30\%.
  \item Validation on the Hillstrom email-marketing RCT
    ($N = 42{,}613$) and covariate-stratified $\tau(x)$ credible
    band coverage across covariate quintiles
    (Section~\ref{sec:experiments:cate_coverage}), substantiating
    calibration for heterogeneous effects beyond scalar ATE summaries.
  \item Six mapped empirical boundaries
    (Section~\ref{sec:discussion}), each with a pointer to the
    benchmark that produced it.
\end{enumerate}

Code, 17+ reproducible benchmarks, and this paper's source are
released under an open-source licence; the repository URL is omitted
from this submission and will be provided upon acceptance.
